\section{Controlled Experimental Design}
\label{sec:design}

\subsection{Two Architecture-Matched Tracks}

Table~\ref{tab:arms} enumerates the eight core arms. Arms 1--4 share a
ViT-B/16 backbone; Arms 5--8 share ConvNeXt-V2-Base. Within each track the
architecture and parameter count are invariant, and only the initialization
changes. The approximate pretrained-backbone sizes are 86M and 89M
parameters; after attaching the family adapter and common detector head, the
implemented models contain 89,996,801 and 93,988,865 trainable parameters,
respectively. All pretrained checkpoints transfer the encoder only, with the
source classification or reconstruction head removed. Random floors use the
identical target architecture without transferred values.

\begin{table}[t]
\caption{Controlled core arms. All non-random entries are downloaded encoders;
task-specific heads are discarded. Parameter counts describe the complete
detector and are constant within each track.}
\label{tab:arms}
\centering
\scriptsize
\setlength{\tabcolsep}{3.2pt}
\begin{tabularx}{\textwidth}{@{}cllXl@{}}
\toprule
Arm & Track & Role & Released initialization and source recipe & License \\
\midrule
1 & ViT & Random & ViT-B/16, fresh initialization & -- \\
2 & ViT & Optical & SatDINO ViT-B/16; DINO on fMoW-RGB & Apache-2.0 \\
3 & ViT & SAR & SARMAE ViT-B; masked autoencoding on SAR-1M & CC BY-NC 4.0 \\
4 & ViT & ImageNet & AugReg ViT-B/16; supervised ImageNet-1K & Apache-2.0 \\
\midrule
5 & CNN & Random & ConvNeXt-V2-Base, fresh initialization & -- \\
6 & CNN & Optical & reBEN BigEarthNet-v2 Sentinel-2 model & model card \\
7 & CNN & SAR & reBEN BigEarthNet-v2 Sentinel-1 model & model card \\
8 & CNN & ImageNet & FCMAE then supervised ImageNet-1K & Apache-2.0 \\
\bottomrule
\end{tabularx}
\vspace{2pt}
\parbox{\textwidth}{\scriptsize Complete detector sizes are 89,996,801
parameters for ViT and 93,988,865 for CNN. The corresponding backbone sizes
are approximately 86M and 89M, respectively. Exact checkpoint paths, hashes,
and source/license notes are archived with the campaign manifest.}
\end{table}


The four roles are:

\begin{description}
  \item[Random floor.] Standard random initialization for ViT-B/16 or
    ConvNeXt-V2-Base.
  \item[ImageNet.] Supervised AugReg ImageNet-1K for ViT-B/16, and
    FCMAE followed by supervised ImageNet-1K fine-tuning for
    ConvNeXt-V2-Base.
  \item[Optical remote sensing.] SatDINO ViT-B/16 pretrained with DINO on
    fMoW-RGB, and BigEarthNet Sentinel-2 ConvNeXt-V2-Base.
  \item[SAR.] SARMAE ViT-B/16 pretrained with masked autoencoding on SAR-1M,
    and BigEarthNet Sentinel-1 ConvNeXt-V2-Base.
\end{description}

The BigEarthNet Sentinel-1/Sentinel-2 pair is the study's cleanest modality
contrast because model, dataset family, and supervised task are shared. The
SatDINO/SARMAE comparison additionally changes DINO versus masked
autoencoding. Likewise, the matched ImageNet role holds source dataset and
final classification supervision constant across tracks but not the complete
training history. The primary domain comparisons are therefore within-track,
and matched-role cross-track comparisons test architectural persistence.

\subsection{Input and Shared Point Detector}

Every arm receives the same three-channel, decibel-domain tensor
\begin{equation}
  \mathbf{x} = [\mathrm{VH},\ \mathrm{VV},\ \mathrm{VH}-\mathrm{VV}].
  \label{eq:channels}
\end{equation}
The channel construction is fixed before the model and is never tuned by arm.
Each released encoder is instantiated through the documented \texttt{timm}
\texttt{in\_chans=3} path~\cite{wightman2019timm}. This leaves already three-channel stems unchanged in
shape and applies the library's repeat-with-rescaling adaptation where the
source checkpoint used a different band count. No hand-written patch-embed or
stem surgery is used.

The detector predicts one vessel-presence heatmap on a common stride-4 grid.
For a $512\times512$ training crop, the output is $128\times128$. Each vessel
point becomes a Gaussian target with standard deviation two output pixels;
LOW-confidence annotations define ignore neighborhoods of radius three output
pixels. Training uses the penalty-reduced focal objective associated with
CenterNet-style point detection~\cite{lin2017focal,zhou2019objects}. The ViT
adapter receives stride-16 features, while the CNN adapter receives stride-32
features and performs one additional learned upsampling stage. Both terminate
in the same 256-channel head and stride-4 heatmap. The different pre-head
resolutions are an inherent architecture-family distinction and are disclosed
rather than tuned away.

At inference, local peaks above a permissive candidate floor of 0.05 are
mapped to scene coordinates. Distance non-maximum suppression uses 120~m, and
the operating threshold is selected only on the designated development
scenes as described in Sect.~\ref{sec:data-evaluation}. Whole scenes are tiled
with $512\times512$ windows at stride 384 and predictions are reconciled in
geographic coordinates. Figure~\ref{fig:detector} summarizes the shared path
from input to matched detections.

\begin{figure}[t]
\centering
\resizebox{\textwidth}{!}{%
\begin{tikzpicture}[font=\sffamily\scriptsize, line width=0.7pt,
  >={Stealth[length=2.2mm, width=1.8mm]},
  stage/.style={draw, rounded corners=2pt, align=center, inner sep=3pt,
                minimum height=8.5mm, text=figink},
  data/.style={stage, fill=black!6, draw=black!45},
  enc/.style={stage, fill=figblue!12, draw=figblue!70!black, minimum width=21mm},
  neck/.style={stage, fill=figblue!6, draw=figblue!50!black},
  cal/.style={stage, fill=figorange!12, draw=figorange!75!black},
  endst/.style={stage, fill=figgreen!10, draw=figgreen!70!black},
  flow/.style={->, draw=figink, rounded corners=1.5pt}]
\node[data] (inp) {$[\mathrm{VH},\mathrm{VV},\mathrm{VH{-}VV}]$\\ $512^2$ crop};
\node[enc] (vit) at ($(inp.east)+(14mm,7.5mm)$) [anchor=west]
  {ViT-B/16\\ stride 16};
\node[enc] (cnn) at ($(inp.east)+(14mm,-7.5mm)$) [anchor=west]
  {ConvNeXt-V2-B\\ stride 32, $+$up};
\begin{scope}[on background layer]
\node[draw=black!35, dashed, rounded corners=3pt, inner sep=2.5mm,
      fit=(vit)(cnn), fill=black!2] (encbox) {};
\end{scope}
\node[anchor=south, font=\sffamily\tiny, text=black!60]
  at (encbox.north) {encoder: one per track};
\node[neck] (head) at ($(encbox.east)+(9mm,0)$) [anchor=west]
  {256-ch head\\ $128^2$ heatmap};
\node[neck, right=5mm of head] (dec) {peaks $\geq 0.05$\\ 120 m NMS};
\node[cal, right=5mm of dec] (op) {$\tau_{\mathrm{dev}}$\\ dev-selected};
\node[endst, right=5mm of op] (mat) {200 m\\ geo match};
\draw[flow] (inp.east) -- ($(inp.east)+(5mm,0)$) |- (vit.west);
\draw[flow] (inp.east) -- ($(inp.east)+(5mm,0)$) |- (cnn.west);
\draw[flow] (encbox.east) -- (head.west)
  node[midway, above, font=\sffamily\tiny, text=black!60] {};
\draw[flow] (head) -- (dec);
\draw[flow] (dec) -- (op);
\draw[flow] (op) -- (mat);
\end{tikzpicture}}
\caption{The shared detection path: only the encoder differs by track; one
head, decode rule, dev-selected threshold, and frozen scorer serve every arm.}
\label{fig:detector}
\end{figure}


\subsection{Uniform Fine-Tuning Recipe}

All eight arms use one PyTorch Lightning module, one data module, and one
trainer configuration~\cite{falcon2019lightning}. AdamW~\cite{loshchilov2019decoupled} uses learning rate
$10^{-4}$, weight decay 0.05, layer-wise decay 0.65, and gradient clipping at
1.0. The schedule has five warm-up epochs followed by cosine decay, with a
maximum of 50 epochs. Full-scene development evaluation occurs every five
epochs, and early stopping uses patience four development evaluations. A
foreground-balanced sampler targets a 0.5 foreground fraction. Shared
augmentations comprise horizontal and vertical flips, $90^{\circ}$ rotations,
a vessel-biased crop policy for 70\% of crops, and a common intensity jitter
of $\pm1$~dB. No hyperparameter is selected per arm or label fraction.

The core matrix uses micro-batch 16, accumulation 1, and therefore effective
batch 16 for every cell. Each run fine-tunes the full encoder and detector
using Lightning \texttt{32-true}. The execution platform is one V100 GPU per
process with no distributed data parallelism. Both CUDA matrix-multiply
and cuDNN TensorFloat-32 paths are disabled before model initialization and
asserted again in training and inference subprocesses. Hardware identifiers,
backend state, code and environment hashes, checkpoint hashes, parameter
counts, and measured runtime are recorded with each cell.

Randomness is fixed to seed 0 for Python, NumPy, PyTorch, and data-loader
workers. The design consequently contains $8\times4=32$ core fine-tunes, not
seed replications. Results are reported as seed-0 point estimates without
error bars or empirical seed-noise claims.

\subsection{Nested Label Fractions and Registered Contrasts}

The four nominal fractions correspond to exact nested training sets of 12,
28, 56, and 111 scenes. If $S_{10}$, $S_{25}$, $S_{50}$, and $S_{100}$ denote
these sets, then
\begin{equation}
  S_{10}\subset S_{25}\subset S_{50}\subset S_{100}.
  \label{eq:nesting}
\end{equation}
The same sets are used by all arms, so a fraction never changes which scenes
one method sees relative to another.

For track $t$, role $r$, and fraction $f$, let $F^{(t)}_{r,f}$ denote
development F1 at that cell's checkpoint-bound operating threshold. In
addition to absolute F1, the analysis reports transfer over the corresponding
same-track random floor and the within-track SAR--optical contrast at every
registered fraction (Sect.~\ref{sec:transfer-results}). A non-monotone curve is reported
as observed rather than smoothed or clamped.

Two external systems, a COCO-pretrained YOLO26 detector and zero-shot
LocateAnything-3B, provide
context~\cite{jocher2026yolo26,wang2026locateanything};
Sect.~\ref{sec:final-results} reports them separately from the controlled
curves.
